\section{FUTURE WORK}

This section will provide practical steps a VERACISM might take post initial controlled deployment to reinforce VERACISM. The planned activities will concentrate on the following areas of resilience, security governance, observability, performance and standards alignment. It will also adhere to the same Assessment criteria established in the Project Assessment baseline to assess for the changes of the project in the future.

\subsection{Reliability and Resilience}

A retry policy with exponential backoff for transient communication failures between Lighthouse will be added, improving the scan and storage flow. Additionally, circuit breaker logic will be introduced to avoid repeated calls to a unhealthy upstream service. A queued upload workflow will further enhance the reliability by temporarily buffering and replay the upload when external services recover.

\subsection{Functional and Integration Coverage Expansion}

Future cycles should build the automated and manual coverage in the same manner in which test cases from the Project Assessment chapter were planned. Functional checking with UI will be enhanced with deeper UI-driven functional checks, integration check for service to service failures (wider), Evidence capture for pass, partial and fail (additional).

\subsection{Audit and Verification Logging}

The verification should continue saving the log entries to the VerificationLog table to track all verification operations. The tenant scoping applied to audit views will also be strictly enforced in API responses and the UI filters. The work will help to reinforce accountability and institutional compliance reviews.

\subsection{Access Control Architecture}

Current tenant ownership checks are effective, but they are distributed across
several controllers. Future work should centralize authorization logic through
a shared authorization layer to reduce duplication and lower the risk of
missing ownership checks in future endpoints. Candidate approaches include
policy-based handlers, reusable ownership validation services, and centralized
authorization middleware for tenant-scoped resources.

Centralising it will also help you have consistent security review in the future. Instead of reviewing each individual controller, reviewers will review a common authorization path, and ensure that new modules conform to the same access control rules.

\subsection{Performance and Capacity Expansion}

Future performance studies will cover more end-to-end points than just public verification, particularly the complete load test for both upload and scanning tests while doing both tests concurrently. Measuring traffic with pilot institutions, then tuning rate-limit thresholds in the reverse proxy level will be done. Endurance tests with long durations will also be added to test for maintaining memory stability and constant throughput.

To ensure security and scanning pipelines are not overwhelmed with excess or abuse of uploads per tenant, a per-tenant upload limit will also be enforced.  This control will keep one tenant from using all the shared scanning resources and aid in maintaining equal distribution of platform usage among institutions. Rate limit will depend on the identity of the tenant, Bulk volume upload, Queue capacity in scan queue on the system, and seen traffic patterns during pilot implementation.

The baseline measures as described in the Project Assessment chapter will be carried forward in the next performance cycle for consistency across versions and environments in terms of throughput, percentile latency, and error rate.

\subsection{Security Hardening and Monitoring}

The platform will be further enhanced with periodic dependency review, continuous security regression testing and automated notifications for unusual access patterns. Scan results, authorization events, upload activity and service health indicators will be collated in a centralized security dashboard to provide faster incident response.

For future testing, the set background principles of automated testing and manual testing will be retained, but with increased linkage of results to evidence of remediation. The performance and security requirements determined in the Project Assessment chapter will also be kept as a reference for additional platforms to be added in the future. This will keep new features from weakening previously tested controls or excessively slowing the response time, or introducing an escalation of risk in the upload/verification/scan process.

\subsection{User Acceptance and Deployment Readiness}

The role-based task scripts for administrators, tenant users, integrators, and external verifiers from the Project Assessment chapter will be formalized as user acceptance testing. This will help in ensuring institutional onboarding and operational handover with more assurance of deployment readiness evidence.

\subsection{Interoperability and Research Extensions}

In the future, this will be investigated with interoperability with verifiable credential frameworks and education metadata standards for the exchange of documents between institutions. Other research aims to investigate verification methods for selective disclosure of student records that preserve privacy.

\subsection{Deployment and Adoption Roadmap}

There will be a rollout at partner institutions ensuring the product validation of governance, onboarding effort, and operational support needs in a staged manner. Samples of user training materials, playbooks for tenant help desk, and tenant onboarding templates will be made more user-friendly and decrease implementation friction.

